% !TEX program = xelatex
\documentclass[cn, chinesefont=nofont, color=black, device=normal, 11pt]{elegantnote}

% 中文字体设置 (遵循 skill 规范)
\setCJKmainfont[AutoFakeSlant=0.2]{Noto Serif CJK SC}
\setCJKsansfont{Noto Sans CJK SC}
\setCJKmonofont{Noto Sans CJK SC}

\title{Vision\_Board 研发日志}
\author{Wenbo}
\date{2026-04-12}

\begin{document}
\maketitle

\section{今日目标}
\begin{itemize}
    \item 重新梳理基于 micro-ROS 与双通道通信的视觉传感器架构规划。
    \item 输出全英文重构的 SPSC 架构以及项目里程碑 Roadmap。
    \item 搭建项目的独立 \LaTeX{} 每日日志系统，配置 Git 防止过程文件污染。
\end{itemize}

\section{技术梳理与难点记录}
今日彻底明确了工程在未来的几个技术硬性指标（即从玩具代码走向工程化）：
\begin{enumerate}
    \item \textbf{内存架构管理}：彻底禁用了 \texttt{new/malloc/std::vector}，拥抱纯静态 SPSC 内存池。修复了由 \texttt{hal\_entry} 函数盲目返回导致主线程死亡、进而拖垮 \texttt{tidle0} 空闲线程清扫器引发的 Stack Overflow 崩溃事故。
    \item \textbf{网络异构传输}：确立了高频 TCP/HTTP（流媒体视频推流）加 低频 UDP/micro-ROS（语义坐标发布）的双通道方案。
    \item \textbf{并发模型}：全系统运用非阻塞模式配合 \texttt{rt\_sem\_take} 休眠唤醒工作流。
\end{enumerate}

\section{问题与反思}
\begin{itemize}
    \item 发现 \texttt{clangd} 对于部分底层编译器私有宏（如 \texttt{ALIGN}）解析时存在跨行牵连报错问题，后均通过遵循现代 C++11 \texttt{alignas} 和标准 GNU 属性 \texttt{\_\_attribute\_\_((aligned(x)))} 规避。
\end{itemize}

\section{明日计划}
\begin{itemize}
    \item 继续完善 micro-ROS 基础通信环境。
    \item 实现轻量级静态内存图像缓冲区。
\end{itemize}

\end{document}
